\documentclass[11pt]{article}
\usepackage[T1]{fontenc}
\usepackage[utf8]{inputenc}
\usepackage[margin=1in]{geometry}
\usepackage{amsmath,amssymb}
\usepackage{microtype}
\usepackage[hidelinks]{hyperref}
\newcommand{\Pb}{\mathbb P}
\newcommand{\E}{\mathbb E}
\DeclareMathOperator{\Beta}{Beta}
\DeclareMathOperator{\BR}{BR}
\allowdisplaybreaks[1]
\title{Competition with a Fixed Outside Option}
\author{Progress note --- transcription of handwritten notes}
\date{\today}

\begin{document}
\maketitle
In this note we present a competition model where the customer chooses
between two sellers and a known outside option of fixed quality $p_0$.

\section{Elements of the model}

\begin{itemize}
\item Two symmetric providers, indexed by $j\in\{A,B\}$.
\item Two products, indexed by $i\in\{1,2\}$.
\item Success probabilities: $0<p_1<p_2<1$.
\item Costs: $c_1<c_2$.
\item Revenue: $R>c_2$.
\item Outside option quality $p_0\in(0,1)$, fixed and known to all agents.
\item Discount factor $\gamma\in(0,1)$ per calendar period.
\end{itemize}

\section{Customer beliefs.}
Recall that in the forgetting model, past observations are discounted by a factor $\rho$ whenever a new outcome is observed. Thus, after $k$ interactions with the seller, the total effective information satisfies
$$
S_k+F_k=\rho(S_{k-1}+F_{k-1})+1,
$$
with \(S_0+F_0=0\). Iterating this expression yields
$$
S_k+F_k=1+\rho+\cdots+\rho^{k-1}
=\frac{1-\rho^k}{1-\rho} \quad \to 1/(1-\rho)\quad\text{as }k\to\infty.
$$

To study the limiting behavior, we normalize the state as 
$$x_s=(1-\rho)S_k$$ 
$$x_f=(1-\rho)F_k$$

Since $x_s+x_f=1-\rho^k\to1$, we focus on the limiting boundary $x_s+x_f=1$, which can be represented by a single variable $x=x_s$, with $x_f=1-x$. Then $x\in [0,1]$ when $\rho \in [1/2,1)$. 

For each seller we will call this normalized state variable $x_j$ for $j\in\{A,B\}$ that will represent the proportion of successes for provider $j$ in the long run. The customer observes the state
$$
x=(x_A,x_B),
$$
Her beliefs are
\begin{align}
\theta_A\mid x_A&\sim\Beta(1+\frac{x_A}{1-\rho},1+\frac{1-x_A}{1-\rho}),\\
\theta_B\mid x_B&\sim\Beta(1+\frac{x_B}{1-\rho},1+\frac{1-x_B}{1-\rho}).
\end{align}


\paragraph{Customer choice.}
The customer uses Thompson sampling (TS), drawing independently from the
two belief distributions and choosing
\begin{equation}
y_t=
\begin{cases}
A,&\widetilde\theta_A>\max\{\widetilde\theta_B,p_0\},\\
B,&\widetilde\theta_B>\max\{\widetilde\theta_A,p_0\},\\
0,&p_0>\max\{\widetilde\theta_A,\widetilde\theta_B\},
\end{cases}
\end{equation}
where $0$ denotes the outside option. Ties have probability zero.
The selection probabilities are
\begin{align}
D_A(x_A,x_B;p_0)
&=\Pb(\widetilde\theta_A>\max\{\widetilde\theta_B,p_0\}\mid x_A,x_B),\\
D_B(x_A,x_B;p_0)
&=\Pb(\widetilde\theta_B>\max\{\widetilde\theta_A,p_0\}\mid x_A,x_B),\\
D_0(x_A,x_B;p_0)
&=\Pb(\widetilde\theta_A<p_0\mid x_A)
  \Pb(\widetilde\theta_B<p_0\mid x_B).
\end{align}
Hence $D_A+D_B+D_0=1$.

\section{Timing, state transitions, and strategies}

\paragraph{Timing.}
\begin{enumerate}
    \item The customer observes $x=(x_A,x_B)$, while each
    seller $j$ observes only its own reputation $x_j$.
    
    \item The customer chooses seller $A$, seller $B$, or the outside
    option using Thompson sampling.
    
    \item If a seller is selected, it chooses a product $i\in\{1,2\}$
    according to its strategy $\sigma_j(\cdot\mid x_j)$
    and receives profit $R-c_i$.
    
    \item If a seller is selected, the customer observes the outcome,
    but not the product.
    
    \item Only the selected seller's reputation is updated.
    If the outside option is chosen, neither reputation changes and
    both sellers receive zero profit.
\end{enumerate}
Thus, forgetting occurs only when a new outcome is observed from a seller,
whereas discounting applies in every calendar period.

\paragraph{State transitions.}
Superscripts $s$ and $f$ denote success and failure,
respectively. The state transitions are
\begin{equation}
\begin{aligned}
T_A^s(x)&=(\rho x_A+1-\rho,x_B),\\
T_A^f(x)&=(\rho x_A,x_B),\\
T_B^s(x)&=(x_A,\rho x_B+1-\rho),\\
T_B^f(x)&=(x_A,\rho x_B).
\end{aligned}
\end{equation}
For an individual seller, we write these updates as
$T_\rho^s(x_j)=\rho x_j+1-\rho$ and $T_\rho^f(x_j)=\rho x_j$.

\paragraph{Strategies.}
We consider stationary mixed strategies that depend only
on the seller's own reputation. For each seller $j\in\{A,B\}$,
a strategy is a mapping
\begin{equation}
\sigma_j:[0,1]\rightarrow\Delta(\{1,2\}),
\end{equation}
where $\sigma_j(i\mid x_j)$ denotes the probability that
seller $j$ chooses product $i$ when selected at reputation
$x_j$.

\section{Oblivious equilibrium: Distribution}

\paragraph{Oblivious approximation:}
Throughout this section, $p_0$ is fixed. It is a parameter of the model,
rather than a conjecture determined by the equilibrium conditions.
Suppose that each seller $j\in\{A,B\}$ assumes that
its competitor's reputation follows a stationary distribution
$\mu_{-j}$ on $[0,1]$, measured in calendar time before customer selection.

Instead of observing the competitor's current reputation,
seller $j$ computes its expected demand by averaging over
this distribution. Write $D(x_j,y;p_0)=D_A(x_j,y;p_0)$, with the
focal seller's reputation as the first argument. By symmetry, the same
function applies to either seller. Then
\begin{equation}
d_j(x_j;\mu_{-j},p_0)
=
\int_0^1 D(x_j,y;p_0)\,\mu_{-j}(dy).
\end{equation}

Thus, each seller's demand depends only on its own reputation
and the conjectured distribution of its competitor, given $p_0$.
When optimizing, the seller treats this distribution as fixed.

\paragraph{Seller's problem:}
Let $U_j(x_j)$ denote seller $j$'s optimal expected discounted
profit before customer selection. We leave the dependence of $U_j$
on the fixed parameter $p_0$ implicit. Given the conjectured distribution
$\mu_{-j}$, the Bellman equation is
\begin{equation}
\begin{aligned}
U_j(x_j) ={}& \gamma[1-d_j(x_j;\mu_{-j},p_0)]U_j(x_j)\\
&+d_j(x_j;\mu_{-j},p_0)
\max_{i\in\{1,2\}}\Big\{ R-c_i\\
&\qquad+\gamma\left[ p_i U_j(T_\rho^s(x_j)) + (1-p_i)U_j(T_\rho^f(x_j))
\right] \Big\}.
\end{aligned}
\end{equation}

If the seller is not selected, either its competitor or the outside option
is chosen. In both cases it receives no profit and its own reputation
remains unchanged, which explains the term $\gamma(1-d_j)U_j$.
If selected, it chooses a product to maximize its current profit and
expected future value.

\paragraph{Stationary distribution:}
Given a strategy $\sigma_j$, the probability of success conditional on
seller $j$ being selected is
\begin{equation}
p_{\sigma_j}(x_j) = \sigma_j(1\mid x_j)p_1 + \sigma_j(2\mid x_j)p_2.
\end{equation}

Given a strategy profile $\sigma=(\sigma_A,\sigma_B)$,
let $P^{\sigma,p_0}(x,C)$ denote the probability that the
joint reputation state belongs to a measurable set
$C\subseteq[0,1]^2$ in the next period, given the
current state $x=(x_A,x_B)$. The transition probabilities are
\begin{equation}
\begin{aligned}
P^{\sigma,p_0}(x,C)
={}&D_0(x_A,x_B;p_0)\mathbf{1}\{x\in C\}\\
&+\sum_{j\in\{A,B\}}D_j(x_A,x_B;p_0)
\Big[p_{\sigma_j}(x_j)\mathbf{1}\{T_j^s(x)\in C\}\\
&\hspace{42mm}+(1-p_{\sigma_j}(x_j))
\mathbf{1}\{T_j^f(x)\in C\}\Big].
\end{aligned}
\end{equation}
The first term accounts for the outside option: both reputations remain
unchanged. In particular, the probability that $B$ sells is $D_B$, rather
than $1-D_A$.

\paragraph{Oblivious equilibrium:}
For the given $p_0$, a symmetric oblivious equilibrium is a mixed strategy $\sigma^*$,
a reputation distribution $\mu^*$, and a joint stationary
distribution $\nu^*$ satisfying the following conditions.
Here $U_{\mu^*}$ is the value from the seller's problem with
$\mu_{-j}=\mu^*$ and the given $p_0$.

\textit{Optimality:}
\begin{equation}
\begin{aligned}
\operatorname{supp}\sigma^*(\cdot\mid x)
\subseteq
\arg\max_{i\in\{1,2\}}
\Big\{
R-c_i
+\gamma\big[
&p_iU_{\mu^*}(T_\rho^s(x))\\
&+(1-p_i)U_{\mu^*}(T_\rho^f(x))
\big]
\Big\}.
\end{aligned}
\end{equation}

\textit{Consistency:}
\begin{equation}
\begin{aligned}
\nu^*&=\nu^*P^{\sigma^*,p_0},\\
\mu^*(C)&=\nu^*(C\times[0,1])
=\nu^*([0,1]\times C),
\qquad \forall C\subseteq[0,1]
\text{ measurable}.
\end{aligned}
\end{equation}

Here $P^{\sigma^*,p_0}$ denotes the joint transition kernel
of the actual two-seller game when both sellers follow
$\sigma^*$ and face the same fixed outside option $p_0$.

The first condition requires each seller to choose only
optimal products given the conjectured distribution $\mu^*$.
The second requires this distribution to coincide with the
marginal distributions generated when both sellers follow
$\sigma^*$ in the actual game, including periods in which the outside
option is chosen. No independence between the two reputations is assumed.

\section{Product choice as the number of sellers grows}

Consider $N$ sellers and one customer choice per calendar period,
keeping $p_0$, $\rho<1$,
$\gamma<1$, and all other parameters fixed. Write
$D_j^{(N)}(x;p_0)$ for seller $j$'s selection probability at the joint
reputation state $x=(x_1,\ldots,x_N)$, and $U_N(x_j)$ for its oblivious
value in the same Bellman equation as above.

\paragraph{Proposition.}
There exists a finite $N_0$ such that, for every $N\geq N_0$, product~1
is strictly optimal at every reputation, for any conjectured joint
distribution of rival reputations. In particular, every symmetric
distribution oblivious equilibrium satisfies
$\sigma^*(1\mid x_j)=1$ for all $x_j\in[0,1]$ when $N\geq N_0$.

\paragraph{Proof.}
Consider that the customer selects seller $j$ when its TS draw exceeds
both $p_0$ and every rival's draw. Ignoring $p_0$ can only increase this
probability, so
\begin{equation}
\begin{aligned}
D_j^{(N)}(x;p_0)
&=\Pb\!\left(
\widetilde\theta_j>\max\{p_0,\max_{k\neq j}\widetilde\theta_k\}
\mid x\right)\\
&\leq\Pb(\widetilde\theta_j>\widetilde\theta_k
\text{ for all }k\neq j\mid x).
\end{aligned}
\end{equation}
Now, write $\widetilde\theta_j(z)$ for seller $j$'s TS draw at reputation
$z$. If $x'>x$, then
\begingroup
\renewcommand{\theHequation}{stochastic-larger-draw}
\begin{equation}
\Pb(\widetilde\theta_j(x')>t)
\geq\Pb(\widetilde\theta_j(x)>t),
\qquad t\in(0,1).
\tag{*}\label{eq:stochastic-larger-draw}
\end{equation}
\endgroup
We prove \eqref{eq:stochastic-larger-draw} below. The most favorable case
for seller $j$ is therefore to have reputation $1$, while every rival
has reputation $0$. Hence
\begin{equation}
\Pb(\widetilde\theta_j>\widetilde\theta_k
\text{ for all }k\neq j\mid x)
\leq\Pb(\widetilde\theta_j(1)>\widetilde\theta_k(0)
\text{ for all }k\neq j).
\end{equation}

Consider that $\widetilde\theta_j(1)\sim\Beta(a,1)$ and
$\widetilde\theta_k(0)\sim\Beta(1,a)$, where $a=1+1/(1-\rho)$.
For $t\in(0,1)$,
\begin{equation}
\Pb(\widetilde\theta_k(0)<t)
=\int_0^t a(1-s)^{a-1}\,ds
=1-(1-t)^a.
\end{equation}
Since the TS draws are independent conditional on reputations, averaging
over seller $j$'s draw gives
\begin{equation}
\begin{aligned}
&\Pb(\widetilde\theta_j(1)>\widetilde\theta_k(0)
\text{ for all }k\neq j)\\
&\qquad=\int_0^1
\big[\Pb(\widetilde\theta_k(0)<t)\big]^{N-1}
a t^{a-1}\,dt\\
&\qquad=\int_0^1\big[1-(1-t)^a\big]^{N-1}a t^{a-1}\,dt\\
&\qquad=\E\!\left[
\big(1-(1-\widetilde\theta_j(1))^a\big)^{N-1}\right].
\end{aligned}
\end{equation}
The bound is valid at every joint reputation state, so
\begin{equation}
\sup_{x\in[0,1]^N}D_j^{(N)}(x;p_0)
\leq\E\!\left[
\big(1-(1-\widetilde\theta_j(1))^a\big)^{N-1}\right].
\end{equation}

Since $\widetilde\theta_j(1)\in(0,1)$ almost surely,
\begin{equation}
\big(1-(1-\widetilde\theta_j(1))^a\big)^{N-1}
\longrightarrow0\quad\text{almost surely},
\end{equation}
and this expression lies between $0$ and $1$. By dominated convergence,
\begin{equation}
\begin{aligned}
&\lim_{N\to\infty}\E\!\left[
\big(1-(1-\widetilde\theta_j(1))^a\big)^{N-1}\right]\\
&\qquad=\E\!\left[
\lim_{N\to\infty}
\big(1-(1-\widetilde\theta_j(1))^a\big)^{N-1}\right]=0.
\end{aligned}
\end{equation}
It follows that
\begin{equation}
\sup_{x\in[0,1]^N}D_j^{(N)}(x;p_0)\longrightarrow0.
\end{equation}
The seller's oblivious demand averages $D_j^{(N)}$ over the conjectured
joint distribution of rival reputations, so the same bound applies to
any such distribution, including correlated reputations.

Now recall that profit per sale is bounded by $R-c_1$. Therefore,
\begin{equation}
\begin{aligned}
0\leq U_N(x_j)
&\leq\sum_{t=0}^{\infty}\gamma^t(R-c_1)
\sup_{x\in[0,1]^N}D_j^{(N)}(x;p_0)\\
&=\frac{R-c_1}{1-\gamma}
\sup_{x\in[0,1]^N}D_j^{(N)}(x;p_0)
\longrightarrow0.
\end{aligned}
\end{equation}
Using this in the optimality condition, and writing
$\Delta c=c_2-c_1>0$ and $\Delta p=p_2-p_1>0$, we obtain
\begin{equation}
\begin{aligned}
&\lim_{N\to\infty}\Big\{
-\Delta c+\gamma\Delta p
\big[U_N(T_\rho^s(x_j))-U_N(T_\rho^f(x_j))\big]\Big\}\\
&\qquad=-\Delta c+\gamma\Delta p(0-0)
=-\Delta c<0.
\end{aligned}
\end{equation}
The convergence is uniform in reputation and rival conjectures.
Thus, product~1 is strictly optimal at every reputation for all
sufficiently large $N$.\hfill$\square$

\paragraph{Proof of \eqref{eq:stochastic-larger-draw}.}
Let $f_x$ denote the posterior density at reputation $x$. For $x'>x$,
\begin{equation}
\frac{f_{x'}(t)}{f_x(t)}
=C\left(\frac{t}{1-t}\right)^{(x'-x)/(1-\rho)},
\qquad C>0,
\end{equation}
where $C$ does not depend on $t$. This ratio increases continuously
from $0$ to infinity. Hence the densities cross once, at some
$t_*\in(0,1)$: $f_{x'}\leq f_x$ below $t_*$ and $f_{x'}\geq f_x$ above it.
Since both densities integrate to $1$,
\begin{equation}
\begin{aligned}
\int_0^t\big[f_{x'}(s)-f_x(s)\big]\,ds
&\leq0,\qquad t\leq t_*,\\
\int_0^t\big[f_{x'}(s)-f_x(s)\big]\,ds
&=-\int_t^1\big[f_{x'}(s)-f_x(s)\big]\,ds\leq0,
\quad t\geq t_*.
\end{aligned}
\end{equation}
Taking complements gives
\begin{equation}
\Pb(\widetilde\theta_j(x')>t)
=1-\int_0^t f_{x'}(s)\,ds
\geq1-\int_0^t f_x(s)\,ds
=\Pb(\widetilde\theta_j(x)>t),
\end{equation}
which proves \eqref{eq:stochastic-larger-draw}.\hfill$\square$

\end{document}
